\documentclass[11pt]{article}
\usepackage[margin=1in]{geometry}
\usepackage{amsmath,amssymb}
\usepackage[hidelinks]{hyperref}
\newcommand{\R}{\mathbb R}
\title{Literature status: plasticity of $B_{\ell_1\oplus_p\R}$}
\author{Run \texttt{fa\_banach\_001} --- lane 12}
\date{11 August 2026}

\begin{document}
\maketitle

\noindent\textbf{Status.}
\emph{Literature already answered: the $\ell_1\oplus_p\R$ subcase for every
$1<p<\infty$}.  This packet records an exact later source; it is not a new
proof.

\section{Source question and special case}

Haller, Leo, and Zavarzina's \emph{Two new examples of Banach spaces with a
plastic unit ball}\linebreak (arXiv:2111.10122; published 2022) asks in its introduction,
PDF page 1, whether the unit ball of every Banach space is plastic.  Here a
metric space is plastic when every non-expansive bijection from the space onto
itself is an isometry.  The paper's Theorem 3 proves the particular case
\[
 B_{\ell_1\oplus_2\R}\quad\text{is plastic}.
\]

A natural direct attack is to replace the outer Euclidean norm by an arbitrary
$p$-norm.  The source proof does extend: strict convexity of $\ell_p^2$
replaces the Euclidean antipodality calculation, and the scalar functions
\[
 A_p(t)=1-|t|^p+|1-t|^p,\qquad
 B_p(t)=1-|t|^p+|1+t|^p
\]
are respectively strictly decreasing and strictly increasing on $[-1,1]$.
The remaining layer-rigidity argument is unchanged after replacing squares by
$p$th powers.

\section{Exact later answer}

Nikita Leo's 2024 University of Tartu Master's thesis
\emph{Banachi ruumi \"uhikkera plastilisus} (English title: \emph{Plasticity
of the Unit Ball of a Banach Space}) contains the exact result in Chapter 2.
At the start of the chapter (printed page 14; PDF page 14), the thesis cites
the Haller--Leo--Zavarzina paper, explicitly notes that the paper treats only
$p=2$, and says that the chapter generalizes it to arbitrary real $p>1$.
It then states:
\begin{quote}
\textbf{Theorem 2.1 (translated).} For every $p\in(1,\infty)$, the unit ball
of $\ell_1\oplus_p\R$ is plastic.
\end{quote}
The thesis supplies the proof through the rest of Chapter 2.  Its English
abstract on PDF page 2 also states the all-$p$ result.  The supporting author
therefore knew and stated that this was a generalization of the source paper;
the match is explicit rather than agent-inferred.

\section{Scope and search record}

This settles the entire finite-$p$ family.  The endpoint $p=1$ is isometric to
$\ell_1$ and is already classical.  Neither the thesis identification nor this
packet settles $p=\infty$ or the source's universal question for arbitrary
Banach spaces.

The four run indexes and the local parsed arXiv corpus contained no duplicate.
Bounded exact web searches for plastic $p$-sums first found only the source;
a subsequent primary-source search located the 2024 University of Tartu thesis
and its exact Chapter 2 statement.  The official supporting PDF was inspected
through DSpace, although direct bitstream transfer repeatedly stalled in the
local environment.  Stable repository links and exact locations are preserved
in \texttt{supporting\_thesis\_metadata.md}.

\section*{References}

\begin{enumerate}
\item R. Haller, N. Leo, and O. Zavarzina, \emph{Two new examples of Banach
spaces with a plastic unit ball}, Acta Comment. Univ. Tartu. Math. 26 (2022),
89--101; arXiv:2111.10122.
\item N. Leo, \emph{Banachi ruumi \"uhikkera plastilisus} [\emph{Plasticity
of the Unit Ball of a Banach Space}], Master's thesis, University of Tartu,
2024, Chapter 2, Theorem 2.1, \url{https://hdl.handle.net/10062/100461}.
\end{enumerate}

\end{document}
